\section{Opposing corners divide a side only at its midpoint}
\label{sec:area_law_side_subdivision}

The subdivisions induced by actual opposing cells follow from the
common dyadic origin and neighboring-layer locality; see
\cite[Section~11]{OpenAI2026AreaLaw}.

% Source: 10-geometry.tex, lines 299--306 and 352--356. This proves
% the possible subdivision points, not constancy of the opposing region.

\begin{definition}[The corners of a dyadic cell]
    \label{def:area_law_cell_corner}
    \lean{TNLean.PEPS.AreaLaw.Geometry.dyadicCellCorner}
    \leanok
    \uses{def:area_law_dyadic_cells}
    For origin $o$, natural exponent $\ell$, integer cell index $z$
    and $\varepsilon\in\{0,1\}^2$, the corresponding corner is
    $o+2^\ell(z+\varepsilon)$.
\end{definition}

\begin{theorem}[An opposing corner on a cell side]
    \label{thm:area_law_corner_side_subdivision}
    \lean{TNLean.PEPS.AreaLaw.Geometry.fineLayer_corner_on_side}
    \leanok
    \uses{def:area_law_cell_corner,def:area_law_cell_fan,
        def:area_law_fine_layer_indices}
    Fix a common origin $o$, finite endpoint set $Z$ and integer
    width $C_0\ge2$. Let $z$ and $w$ index actual fine cells of layers
    $k$ and $h$, respectively, with $k\ge50\,000\,000$. Write $[a,b]$
    for any whole side of the cell indexed by $z$. If a corner $x$
    of the cell indexed by $w$ belongs to $[a,b]$, then
    $x=a$, $x=(a+b)/2$, or $x=b$.
\end{theorem}
\begin{proof}\leanok
    \uses{def:area_law_cell_corner,def:area_law_cell_fan,
        thm:area_law_dyadic_closures,thm:area_law_fine_cell_containment,
        thm:area_law_local_layer_indices,thm:area_law_adjacent_fine_meshes,
        def:area_law_affine_mesh,thm:area_law_cell_fan_mark_vertices,
        thm:area_law_belt_marks_closed_cell}
    The shared point belongs to both closed layers, so
    $h\le k+1$ and $k\le h+1$. If $\ell_j$ is the fine exponent
    of layer $j$, it follows that $\ell_k\le\ell_h+1$.
    Put $t=2^{\ell_k}$. The integer coordinates of the opposing
    corner place it on $o+(t/2)\mathbb Z^2$. Along the reference
    side its varying coordinate, measured from $a$ in units of
    $t/2$, is therefore an integer between $0$ and $2$. These
    three integers give the two endpoints and the midpoint.
\end{proof}
